\chapter{Agent Development Frameworks}
\label{sec:agent-development-frameworks}

The transition from a research prototype to a production-grade agent system is one of the most demanding engineering challenges in modern AI development. While academic papers demonstrate impressive capabilities in controlled settings, real-world deployment exposes a host of concerns that go far beyond raw task performance: reliability under adversarial inputs, observability of internal reasoning, testability of complex multi-step workflows, and the operational overhead of serving millions of requests at scale. This section surveys the landscape of agent development frameworks---the tools, libraries, and platforms that have emerged to address these challenges---and provides practical guidance for building, testing, deploying, and iterating on production agent systems.

\section{Motivation: The Engineering Gap}
\label{subsec:engineering-gap}

\begin{keybox}[Why Agent Engineering Is Hard]
Building a capable agent in a Jupyter notebook is straightforward. Building one that runs reliably in production---handling edge cases, recovering from failures, scaling to load, and improving over time---requires a fundamentally different engineering discipline.
\end{keybox}

Research prototypes typically assume a cooperative environment: well-formed inputs, available tools, responsive APIs, and a patient human observer ready to restart the process when something goes wrong. Production agents face none of these luxuries. The engineering gap between prototype and production manifests across several dimensions:

\paragraph{Reliability.}
\label{reliability.}

A production agent must handle tool failures gracefully, recover from partial state corruption, and avoid infinite loops or runaway API calls. Error handling must be systematic, not ad hoc.

\paragraph{Observability.}
\label{observability.}

When an agent produces a wrong answer or takes an unexpected action, operators need to understand \emph{why}. This requires structured logging of every LLM call, tool invocation, and state transition---not just the final output.

\paragraph{Testability.}
\label{testability.}

Agent behavior is non-deterministic and context-dependent, making traditional unit testing insufficient. Comprehensive agent testing requires specialized evaluation harnesses, golden trajectory comparisons, and behavioral test suites.

\paragraph{Deployment.}
\label{deployment.}

Agents are stateful, long-running processes that may span minutes or hours. Serving infrastructure must support async execution, checkpointing, resumption after failures, and multi-tenant isolation.

\paragraph{Iteration.}
\label{iteration.}

Production agents degrade over time as the world changes, APIs evolve, and user behavior shifts. Continuous improvement requires systematic failure analysis, prompt versioning, and fine-tuning pipelines.

\begin{intuitionbox}[The Agent Development Maturity Model]
Agent development follows a maturity progression:

\begin{enumerate}
  \item \textbf{Prototype}: Single-file script, hardcoded prompts, manual testing
  \item \textbf{Alpha}: Modular code, basic error handling, manual evaluation
  \item \textbf{Beta}: Framework-based, automated tests, staging environment
  \item \textbf{Production}: Full observability, CI/CD, auto-scaling, SLAs
  \item \textbf{Mature}: Continuous learning, A/B testing, self-improvement loops
\end{enumerate}

Most teams underestimate the gap between stages 2 and 3.
\end{intuitionbox}

\section{The Agent Development Lifecycle}
\label{subsec:agent-lifecycle}

A structured development lifecycle helps teams move systematically from concept to production. Figure~\ref{fig:agent-lifecycle} illustrates the five major phases.

\begin{figure}[ht!]
\centering
\includegraphics[width=\textwidth]{figures/fig_072_agent-lifecycle.png}
\caption{The agent development lifecycle. Feedback loops at every stage ensure continuous improvement.}
\label{fig:agent-lifecycle}
\end{figure}

\subsection{Phase 1: Design}
\label{phase-1-design}

The design phase establishes the agent’s \emph{capability envelope}---what it can and cannot do---before a single line of code is written.

\textbf{Defining capabilities.} Start with a capability matrix: a structured list of tasks the agent should handle, edge cases it must reject, and behaviors that are explicitly out of scope. This document becomes the basis for evaluation criteria.

\textbf{Tool selection.} Each tool should have a clear purpose, well-defined inputs and outputs, and a failure mode specification. Over-tooling is a common mistake: agents with too many tools suffer from tool selection confusion and increased latency.

\textbf{Constraint specification.} Production agents require explicit constraints: maximum number of tool calls per request, allowed domains for web browsing, data access permissions, and output format requirements. These constraints should be encoded in the system prompt \emph{and} enforced programmatically.

\subsection{Phase 2: Implementation}
\label{phase-2-implementation}

Implementation involves three interleaved concerns: prompt engineering, tool integration, and orchestration logic.

\textbf{Prompt engineering.} System prompts for production agents are living documents that require version control, structured testing, and careful change management. Techniques include chain-of-thought scaffolding, few-shot examples, explicit output format instructions, and persona definition.

\textbf{Tool integration.} Each tool is implemented as a function with a typed interface, comprehensive error handling, and idempotency guarantees where possible. Tool descriptions (used by the LLM to decide when to invoke them) are as important as the tool implementations themselves.

\textbf{Orchestration.} The orchestration layer manages the agent loop: calling the LLM, parsing tool calls, executing tools, updating state, and deciding when to terminate. Framework choice (Section~\ref{subsec:frameworks}) significantly impacts how this layer is structured.

\subsection{Phase 3: Testing}
\label{phase-3-testing}

Agent testing is covered in depth in Section~\ref{subsec:agent-testing}. The key principle is \emph{test at multiple granularities}: individual tools, complete agent loops, and end-to-end user scenarios.

\subsection{Phase 4: Deployment}
\label{phase-4-deployment}

Deployment concerns are covered in Section~\ref{subsec:production-deployment}. Key decisions include synchronous vs.~asynchronous execution, state persistence strategy, and scaling architecture.

\subsection{Phase 5: Iteration}
\label{phase-5-iteration}

The iteration phase closes the loop between production behavior and system improvement. It requires:

\begin{itemize}
  \item \textbf{Failure logging}: Every agent failure is logged with full context (input, trajectory, error)
  \item \textbf{Failure categorization}: Failures are classified by type (tool error, reasoning error, hallucination, loop) to identify systemic issues
  \item \textbf{Prompt updates}: Prompt changes are tested against regression suites before deployment
  \item \textbf{Fine-tuning}: When prompt engineering reaches its limits, fine-tuning on curated trajectories can improve performance
  \item \textbf{A/B testing}: New agent versions are tested against production traffic with statistical rigor
\end{itemize}

\section{Major Frameworks: A Deep Dive}
\label{subsec:frameworks}

The agent framework ecosystem has grown rapidly, with each framework reflecting different design philosophies and target use cases. We examine the most widely adopted frameworks in depth.

\subsection{LangGraph}
\label{subsubsec:langgraph}

LangGraph~\cite{langchain2024langgraph}, developed by LangChain Inc., models agent execution as a \emph{directed graph} where nodes represent computation steps and edges represent transitions between steps. This graph-based abstraction provides explicit control over agent flow, making it easier to reason about, test, and debug complex multi-step behaviors.

\paragraph{Core Concepts.}
\label{core-concepts.}

\begin{itemize}
  \item \textbf{State}: A typed dictionary (using Python’s \texttt{TypedDict} or Pydantic) that flows through the graph and is updated by each node
  \item \textbf{Nodes}: Python functions that receive the current state and return state updates
  \item \textbf{Edges}: Transitions between nodes, which can be unconditional or conditional (routing based on state)
  \item \textbf{Checkpointing}: Built-in persistence of graph state, enabling pause/resume and human-in-the-loop workflows
  \item \textbf{Subgraphs}: Composable graph components that can be nested within larger graphs
\end{itemize}

\paragraph{State Management.}
\label{state-management.}

LangGraph’s state management is one of its most powerful features. The state schema acts as a contract between nodes, making data flow explicit and type-safe:

\begin{lstlisting}[style=pythonstyle, caption={LangGraph state schema definition}]
from typing import TypedDict, Annotated, List
from langgraph.graph.message import add_messages

class AgentState(TypedDict):
    # Messages accumulate via the add_messages reducer
    messages: Annotated[List[BaseMessage], add_messages]
    # Simple fields are overwritten on each update
    current_tool: str | None
    iteration_count: int
    final_answer: str | None
    error: str | None
\end{lstlisting}

\paragraph{Checkpointing and Human-in-the-Loop.}
\label{checkpointing-and-human-in-the-loop.}

LangGraph’s checkpointer saves graph state after every node execution. This enables:

\begin{itemize}
  \item \textbf{Resumption}: Long-running agents can be paused and resumed without losing progress
  \item \textbf{Human approval}: The graph can pause at designated nodes and wait for human input before proceeding
  \item \textbf{Time travel}: Operators can replay execution from any checkpoint for debugging
\end{itemize}

\begin{lstlisting}[style=pythonstyle, caption={LangGraph checkpointing and human-in-the-loop}]
from langgraph.checkpoint.sqlite import SqliteSaver
from langgraph.graph import StateGraph, START, END

# Persistent checkpointer
memory = SqliteSaver.from_conn_string("agent_state.db")

# Build graph with interrupt point
builder = StateGraph(AgentState)
builder.add_node("plan", plan_node)
builder.add_node("human_review", human_review_node)
builder.add_node("execute", execute_node)

builder.add_edge(START, "plan")
builder.add_edge("plan", "human_review")
builder.add_edge("human_review", "execute")
builder.add_edge("execute", END)

# Compile with checkpointer and interrupt before human_review
graph = builder.compile(
    checkpointer=memory,
    interrupt_before=["human_review"]
)

# Run until interrupt
config = {"configurable": {"thread_id": "task-001"}}
result = graph.invoke({"messages": [HumanMessage("Analyze Q3 sales")]}, config)

# Resume after human provides input
graph.update_state(config, {"human_feedback": "Approved, proceed"})
result = graph.invoke(None, config)  # Resume from checkpoint
\end{lstlisting}

The following two listings combine every element above---state schemas, tool nodes, conditional routing, checkpointing, and invocation---into a complete research agent that iteratively gathers information and synthesizes a report.

\begin{lstlisting}[style=pythonstyle, caption={Research agent -- state, tools, and node functions}]
from typing import TypedDict, Annotated, List
from langchain_openai import ChatOpenAI
from langchain_core.tools import tool
from langchain_core.messages import BaseMessage, HumanMessage, AIMessage
from langgraph.graph import StateGraph, START, END
from langgraph.prebuilt import ToolNode
from langgraph.graph.message import add_messages
from langgraph.checkpoint.sqlite import SqliteSaver

# --- Tool Definitions ---
@tool
def search_web(query: str) -> str:
    """Search the web for current information on a topic."""
    return f"Search results for: {query}"  # stub; call real API

@tool
def read_document(url: str) -> str:
    """Fetch and read the content of a document at a URL."""
    return f"Document content from: {url}"

tools = [search_web, read_document]

# --- State Schema ---
class ResearchState(TypedDict):
    messages: Annotated[List[BaseMessage], add_messages]
    research_topic: str
    iteration: int
    status: str  # "researching" | "drafting" | "done" | "error"

# --- Node Functions ---
def research_node(state: ResearchState) -> dict:
    """LLM decides what to search next or signals completion."""
    llm = ChatOpenAI(model="gpt-4o").bind_tools(tools)
    response = llm.invoke(state["messages"])
    return {"messages": [response], "iteration": state["iteration"] + 1}

def should_continue(state: ResearchState) -> str:
    """Route: tool calls -> execute tools; no calls -> synthesize."""
    last = state["messages"][-1]
    if hasattr(last, "tool_calls") and last.tool_calls:
        return "tools"
    if state["iteration"] >= 10:
        return "error"
    return "synthesize"

def synthesize_node(state: ResearchState) -> dict:
    """Produce final report from accumulated research."""
    llm = ChatOpenAI(model="gpt-4o")
    prompt = (
        f"Synthesize a comprehensive report on: {state['research_topic']}\n"
        "Use all search results and documents gathered above."
    )
    response = llm.invoke(
        state["messages"] + [HumanMessage(content=prompt)]
    )
    return {"messages": [response], "status": "done"}

def error_node(state: ResearchState) -> dict:
    return {"status": "error", "messages": [
        AIMessage(content="Research exceeded maximum iterations.")
    ]}
\end{lstlisting}

\begin{lstlisting}[style=pythonstyle, caption={Research agent -- graph construction and invocation}]
# --- Graph Construction ---
tool_node = ToolNode(tools)
builder = StateGraph(ResearchState)
builder.add_node("research", research_node)
builder.add_node("tools", tool_node)
builder.add_node("synthesize", synthesize_node)
builder.add_node("error", error_node)

builder.add_edge(START, "research")
builder.add_conditional_edges(
    "research", should_continue,
    {"tools": "tools", "synthesize": "synthesize", "error": "error"}
)
builder.add_edge("tools", "research")   # loop back after tool execution
builder.add_edge("synthesize", END)
builder.add_edge("error", END)

# Compile with persistence for conversation memory
with SqliteSaver.from_conn_string(":memory:") as checkpointer:
    graph = builder.compile(checkpointer=checkpointer)

# --- Invoke ---
result = graph.invoke(
    {"messages": [HumanMessage(content="Research recent advances in RLHF")],
     "research_topic": "Recent advances in RLHF",
     "iteration": 0, "status": "researching"},
    config={"configurable": {"thread_id": "research-1"}}
)
\end{lstlisting}

\begin{figure}[ht!]
\centering
\includegraphics[width=0.85\textwidth]{figures/fig_073_langgraph-graph.png}
\caption{LangGraph execution graph for the research agent. Conditional edges implement the tool-use loop and error handling.}
\label{fig:langgraph-graph}
\end{figure}

\subsection{AutoGen (Microsoft)}
\label{subsubsec:autogen}

AutoGen~\cite{wu2023autogen}, developed by Microsoft Research, takes a fundamentally different approach: it models agents as \emph{conversable entities} that communicate through structured message passing. Rather than a single agent loop, AutoGen enables multi-agent conversations where specialized agents collaborate to solve complex tasks.

\paragraph{Conversable Agents.}
\label{conversable-agents.}

Every AutoGen agent is a \texttt{ConversableAgent} with:

\begin{itemize}
  \item A \textbf{system message} defining its role and capabilities
  \item A \textbf{human input mode} controlling when it solicits human input (\texttt{ALWAYS}, \texttt{NEVER}, \texttt{TERMINATE})
  \item A \textbf{code execution config} specifying whether and how it can run code
  \item A \textbf{function map} of callable tools
\end{itemize}

\paragraph{Group Chat Patterns.}
\label{group-chat-patterns.}

AutoGen’s \texttt{GroupChat} enables multiple agents to collaborate in a shared conversation. A \texttt{GroupChatManager} orchestrates turn-taking, either through round-robin, LLM-based speaker selection, or custom routing logic.

\begin{lstlisting}[style=pythonstyle, caption={AutoGen multi-agent group chat}]
import autogen

config_list = [{"model": "gpt-4o", "api_key": os.environ["OPENAI_API_KEY"]}]
llm_config = {"config_list": config_list, "temperature": 0}

# Specialized agents
planner = autogen.AssistantAgent(
    name="Planner",
    system_message="""You are a strategic planner. Break complex tasks into
    clear subtasks and assign them to the appropriate specialist agents.
    Always end your message with a clear action item for another agent.""",
    llm_config=llm_config,
)

coder = autogen.AssistantAgent(
    name="Coder",
    system_message="""You are an expert Python programmer. Write clean,
    well-documented code. Always test your code before presenting it.""",
    llm_config=llm_config,
    code_execution_config={"work_dir": "coding", "use_docker": True},
)

critic = autogen.AssistantAgent(
    name="Critic",
    system_message="""You review code and plans for correctness, efficiency,
    and security. Provide specific, actionable feedback.""",
    llm_config=llm_config,
)

user_proxy = autogen.UserProxyAgent(
    name="UserProxy",
    human_input_mode="TERMINATE",
    max_consecutive_auto_reply=10,
    is_termination_msg=lambda x: "TASK_COMPLETE" in x.get("content", ""),
    code_execution_config={"work_dir": "output", "use_docker": False},
)

# Group chat with LLM-based speaker selection
groupchat = autogen.GroupChat(
    agents=[user_proxy, planner, coder, critic],
    messages=[],
    max_round=20,
    speaker_selection_method="auto",
)
manager = autogen.GroupChatManager(groupchat=groupchat, llm_config=llm_config)

# Initiate the conversation
user_proxy.initiate_chat(
    manager,
    message="Analyze the CSV dataset in 'sales_data.csv' and generate a summary report with visualizations."
)
\end{lstlisting}

\paragraph{Code Execution Agents.}
\label{code-execution-agents.}

AutoGen’s code execution capability is a distinguishing feature. The \texttt{UserProxyAgent} can execute Python and shell code in a sandboxed environment (Docker container or local process), enabling agents to iteratively write, test, and fix code.

\begin{warningbox}[AutoGen Security Considerations]
Code execution agents can run arbitrary code. Always use Docker isolation in production environments. Configure \texttt{code\_execution\_config} with \texttt{"use\_docker": True} and restrict network access. Never run AutoGen code execution agents with elevated privileges.
\end{warningbox}

\subsection{CrewAI}
\label{subsubsec:crewai}

CrewAI~\cite{moura2023crewai} introduces a \emph{role-based} paradigm for multi-agent systems, drawing inspiration from organizational management. Agents are defined by their professional roles, goals, and backstories---a design choice that leverages the LLM’s understanding of human organizational structures.

\paragraph{Core Abstractions.}
\label{core-abstractions.}

\begin{itemize}
  \item \textbf{Agent}: Defined by \texttt{role}, \texttt{goal}, \texttt{backstory}, and available \texttt{tools}
  \item \textbf{Task}: A specific assignment with a \texttt{description}, \texttt{expected\_output}, and assigned \texttt{agent}
  \item \textbf{Crew}: A collection of agents and tasks with an execution \texttt{process} (sequential or hierarchical)
  \item \textbf{Process}: Execution strategy---\texttt{sequential} (tasks run in order) or \texttt{hierarchical} (a manager agent delegates)
\end{itemize}

\begin{lstlisting}[style=pythonstyle, caption={CrewAI role-based agent team}]
from crewai import Agent, Task, Crew, Process
from crewai_tools import SerperDevTool, WebsiteSearchTool

search_tool = SerperDevTool()
web_tool = WebsiteSearchTool()

# Define agents with rich role descriptions
researcher = Agent(
    role="Senior Research Analyst",
    goal="Uncover cutting-edge developments in AI and provide "
         "comprehensive, accurate research summaries",
    backstory="""You are a seasoned research analyst with 15 years of
    experience in technology research. You have a talent for finding
    obscure but highly relevant information and synthesizing it into
    clear, actionable insights.""",
    tools=[search_tool, web_tool],
    verbose=True,
    allow_delegation=False,
)

writer = Agent(
    role="Tech Content Strategist",
    goal="Craft compelling, technically accurate content that "
         "engages both technical and non-technical audiences",
    backstory="""You are a renowned content strategist known for
    translating complex technical concepts into engaging narratives.
    Your writing has appeared in major tech publications.""",
    tools=[web_tool],
    verbose=True,
    allow_delegation=True,
)

# Define tasks with clear expected outputs
research_task = Task(
    description="""Conduct comprehensive research on {topic}.
    Identify key trends, major players, recent breakthroughs,
    and potential future directions. Focus on developments from
    the past 6 months.""",
    expected_output="""A detailed research report with:
    - Executive summary (200 words)
    - Key findings (5-7 bullet points)
    - Detailed analysis (500 words)
    - Sources and citations""",
    agent=researcher,
)

writing_task = Task(
    description="""Using the research provided, write a compelling
    blog post about {topic} for a technical audience.""",
    expected_output="""A polished blog post (800-1000 words) with:
    - Engaging headline
    - Introduction hook
    - 3-4 main sections with subheadings
    - Conclusion with call to action""",
    agent=writer,
    context=[research_task],  # Depends on research output
)

# Assemble the crew
crew = Crew(
    agents=[researcher, writer],
    tasks=[research_task, writing_task],
    process=Process.sequential,
    verbose=2,
)

result = crew.kickoff(inputs={"topic": "Reinforcement Learning for LLMs"})
\end{lstlisting}

\paragraph{Hierarchical Process.}
\label{hierarchical-process.}

In hierarchical mode, CrewAI automatically creates a manager agent that delegates tasks to worker agents based on their roles and capabilities. This mirrors real organizational structures and can handle complex, interdependent workflows without explicit task ordering.

\subsection{OpenAI Assistants API and Agents SDK}
\label{subsubsec:openai-agents}

OpenAI provides two complementary offerings for agent development: the \textbf{Assistants API}, a hosted infrastructure for stateful agents, and the \textbf{Agents SDK}~\cite{openai2024agentssdk} (formerly Swarm), a lightweight Python library for multi-agent orchestration.

\paragraph{Assistants API Architecture.}
\label{assistants-api-architecture.}

The Assistants API manages agent state server-side through three core objects:

\begin{itemize}
  \item \textbf{Assistant}: A configured agent with a model, instructions, and tools
  \item \textbf{Thread}: A persistent conversation history associated with a user session
  \item \textbf{Run}: An execution of an assistant on a thread, with a lifecycle of states (\texttt{queued} $\to$ \texttt{in\_progress} $\to$ \texttt{requires\_action} $\to$ \texttt{completed})
\end{itemize}

\paragraph{Built-in Tools.}
\label{built-in-tools.}

The Assistants API provides three hosted tools that require no external infrastructure:

\begin{itemize}
  \item \textbf{Code Interpreter}: Executes Python in a sandboxed environment with file I/O
  \item \textbf{File Search}: Vector-store-backed retrieval over uploaded documents
  \item \textbf{Web Search}: Real-time web browsing (available in select models)
\end{itemize}

\begin{lstlisting}[style=pythonstyle, caption={OpenAI Assistants API with tool use}]
from openai import OpenAI
import time

client = OpenAI()

# Create a persistent assistant
assistant = client.beta.assistants.create(
    name="Data Analysis Assistant",
    instructions="""You are an expert data analyst. When given data files,
    analyze them thoroughly and provide actionable insights with
    visualizations where appropriate.""",
    model="gpt-4o",
    tools=[
        {"type": "code_interpreter"},
        {"type": "file_search"},
    ],
)

# Create a thread for a user session
thread = client.beta.threads.create()

# Upload a data file
with open("sales_data.csv", "rb") as f:
    file = client.files.create(file=f, purpose="assistants")

# Add a message with the file attachment
client.beta.threads.messages.create(
    thread_id=thread.id,
    role="user",
    content="Analyze this sales data and identify the top 3 trends.",
    attachments=[{"file_id": file.id, "tools": [{"type": "code_interpreter"}]}],
)

# Create and poll a run
run = client.beta.threads.runs.create_and_poll(
    thread_id=thread.id,
    assistant_id=assistant.id,
)

if run.status == "completed":
    messages = client.beta.threads.messages.list(thread_id=thread.id)
    print(messages.data[0].content[0].text.value)
elif run.status == "requires_action":
    # Handle function tool calls
    tool_calls = run.required_action.submit_tool_outputs.tool_calls
    outputs = []
    for tc in tool_calls:
        result = dispatch_tool(tc.function.name, tc.function.arguments)
        outputs.append({"tool_call_id": tc.id, "output": result})
    client.beta.threads.runs.submit_tool_outputs(
        thread_id=thread.id, run_id=run.id, tool_outputs=outputs
    )
\end{lstlisting}

\paragraph{OpenAI Agents SDK: Swarm Patterns.}
\label{openai-agents-sdk-swarm-patterns.}

The Agents SDK provides a lightweight framework for multi-agent handoffs. The key primitive is the \emph{handoff}: an agent can transfer control to another agent, passing along context. This enables modular agent architectures where specialized agents handle specific subtasks.

\begin{lstlisting}[style=pythonstyle, caption={OpenAI Agents SDK with handoffs and guardrails}]
from agents import Agent, Runner, RunConfig, handoff, InputGuardrail, GuardrailFunctionOutput
from pydantic import BaseModel

# Input validation guardrail
class SafetyCheck(BaseModel):
    is_safe: bool
    reason: str

async def safety_guardrail(ctx, agent, input_data):
    result = await Runner.run(
        Agent(
            name="SafetyChecker",
            instructions="Check if the request is safe and appropriate.",
            output_type=SafetyCheck,
        ),
        input_data,
    )
    return GuardrailFunctionOutput(
        output_info=result.final_output,
        tripwire_triggered=not result.final_output.is_safe,
    )

# Specialized agents
billing_agent = Agent(
    name="BillingAgent",
    instructions="Handle billing inquiries, refunds, and payment issues.",
    tools=[lookup_invoice, process_refund],
)

technical_agent = Agent(
    name="TechnicalAgent",
    instructions="Resolve technical issues and bugs.",
    tools=[check_system_status, create_ticket],
)

# Triage agent with handoffs
triage_agent = Agent(
    name="TriageAgent",
    instructions="""Classify customer requests and route to the appropriate
    specialist. Use handoffs to transfer to billing or technical agents.""",
    handoffs=[
        handoff(billing_agent, tool_name_override="transfer_to_billing"),
        handoff(technical_agent, tool_name_override="transfer_to_technical"),
    ],
    input_guardrails=[InputGuardrail(guardrail_function=safety_guardrail)],
)

# Run with tracing enabled
result = await Runner.run(
    triage_agent,
    "I was charged twice for my subscription last month.",
    run_config=RunConfig(tracing_disabled=False),
)
\end{lstlisting}

\subsection{DSPy}
\label{subsubsec:dspy}

DSPy~\cite{khattab2023dspy} (Declarative Self-improving Python) takes a radically different approach to agent development: rather than manually engineering prompts, DSPy \emph{compiles} high-level program specifications into optimized prompts through automated optimization.

\paragraph{Core Philosophy.}
\label{core-philosophy.}

DSPy separates \emph{what} a module should do (its signature) from \emph{how} it should do it (the prompt). Optimizers then search for the best prompts and few-shot examples to maximize a metric on a development set. This makes DSPy programs more robust to model changes and eliminates the need for manual prompt tuning.

\begin{lstlisting}[style=pythonstyle, caption={DSPy signatures and modules}]
import dspy

# Configure the language model
lm = dspy.LM("openai/gpt-4o", temperature=0.0)
dspy.configure(lm=lm)

# Signatures define input/output contracts
class GenerateAnswer(dspy.Signature):
    """Answer questions with factual, concise responses."""
    context: list[str] = dspy.InputField(desc="Relevant passages")
    question: str = dspy.InputField()
    answer: str = dspy.OutputField(desc="Concise factual answer")

class AssessAnswer(dspy.Signature):
    """Assess whether an answer is faithful to the context."""
    context: list[str] = dspy.InputField()
    question: str = dspy.InputField()
    answer: str = dspy.InputField()
    faithful: bool = dspy.OutputField()
    confidence: float = dspy.OutputField(desc="Confidence score 0-1")

# Modules compose signatures into programs
class RAGAgent(dspy.Module):
    def __init__(self, num_passages=3):
        self.retrieve = dspy.Retrieve(k=num_passages)
        self.generate = dspy.ChainOfThought(GenerateAnswer)
        self.assess = dspy.Predict(AssessAnswer)

    def forward(self, question: str) -> dspy.Prediction:
        context = self.retrieve(question).passages
        prediction = self.generate(context=context, question=question)

        # Self-assessment with assertion
        assessment = self.assess(
            context=context,
            question=question,
            answer=prediction.answer,
        )
        dspy.Assert(
            assessment.faithful,
            "Answer not faithful to context "
            "(confidence: " + str(assessment.confidence) + ")"
        )
        return prediction
\end{lstlisting}

\paragraph{Optimizers.}
\label{optimizers.}

DSPy’s optimizers automatically improve program performance:

\begin{lstlisting}[style=pythonstyle, caption={DSPy optimization with MIPRO}]
from dspy.teleprompt import MIPROv2

# Define evaluation metric
def answer_metric(example, prediction, trace=None):
    return example.answer.lower() in prediction.answer.lower()

# Compile with MIPRO optimizer
optimizer = MIPROv2(
    metric=answer_metric,
    auto="medium",  # Controls optimization budget
)

compiled_agent = optimizer.compile(
    RAGAgent(),
    trainset=train_examples,
    num_candidates=30,
    max_bootstrapped_demos=4,
    max_labeled_demos=16,
)

# Save optimized program
compiled_agent.save("optimized_rag_agent.json")
\end{lstlisting}

\begin{intuitionbox}[When to Use DSPy]
DSPy excels when: (1) you have a clear evaluation metric, (2) you have a development dataset of 50+ examples, (3) you need to port your agent across different LLMs, or (4) manual prompt engineering has plateaued. It is less suitable for highly creative tasks where the ``correct'' output is subjective.
\end{intuitionbox}

\subsection{Semantic Kernel (Microsoft)}
\label{subsubsec:semantic-kernel}

Semantic Kernel~\cite{microsoft2023semantickernel} (SK) is Microsoft’s enterprise-focused agent framework, designed for integration with existing software systems and organizational workflows. It provides a \emph{plugin architecture} that allows developers to expose existing business logic as AI-callable functions.

\paragraph{Plugin Architecture.}
\label{plugin-architecture.}

Plugins are collections of functions (``skills'') that the kernel can invoke. They can be defined as:

\begin{itemize}
  \item \textbf{Native functions}: Regular Python/C\# methods decorated with \texttt{@kernel\_function}
  \item \textbf{Prompt functions}: Parameterized prompt templates stored as files
  \item \textbf{OpenAPI plugins}: Auto-generated from OpenAPI specifications
\end{itemize}

\begin{lstlisting}[style=pythonstyle, caption={Semantic Kernel plugin and planner}]
import semantic_kernel as sk
from semantic_kernel.functions import kernel_function
from semantic_kernel.connectors.ai.open_ai import OpenAIChatCompletion

kernel = sk.Kernel()
kernel.add_service(OpenAIChatCompletion(ai_model_id="gpt-4o"))

# Define a native plugin
class EmailPlugin:
    @kernel_function(description="Send an email to a recipient")
    def send_email(self, recipient: str, subject: str, body: str) -> str:
        # Integration with email service
        return f"Email sent to {recipient}: {subject}"

    @kernel_function(description="Search emails by keyword")
    def search_emails(self, query: str, max_results: int = 10) -> str:
        # Integration with email search API
        return f"Found {max_results} emails matching: {query}"

class CalendarPlugin:
    @kernel_function(description="Schedule a meeting")
    def schedule_meeting(
        self, title: str, attendees: str, datetime_str: str
    ) -> str:
        return f"Meeting '{title}' scheduled for {datetime_str}"

# Register plugins
kernel.add_plugin(EmailPlugin(), plugin_name="Email")
kernel.add_plugin(CalendarPlugin(), plugin_name="Calendar")

# Use the function-calling planner
from semantic_kernel.planners import FunctionCallingStepwisePlanner

planner = FunctionCallingStepwisePlanner(service_id="gpt-4o")
result = await planner.invoke(
    kernel,
    "Schedule a meeting with alice@company.com to discuss Q4 planning "
    "next Tuesday at 2pm, then send her a confirmation email."
)
print(str(result))
\end{lstlisting}

\paragraph{Memory and Connectors.}
\label{memory-and-connectors.}

Semantic Kernel’s memory system supports multiple backends (Azure Cognitive Search, Chroma, Pinecone, Weaviate) through a unified interface. The connector system enables integration with enterprise services including Microsoft 365, Azure DevOps, and custom REST APIs.

\paragraph{Enterprise Integration Focus.}
\label{enterprise-integration-focus.}

SK is particularly well-suited for enterprise deployments due to:

\begin{itemize}
  \item Native C\# support for .NET ecosystems
  \item Azure OpenAI integration with managed identity authentication
  \item Compliance-friendly architecture with audit logging
  \item Support for on-premises model deployments
\end{itemize}

\section{Open-Source Agent Tooling}
\label{subsec:open-source-tooling}

Beyond the major commercial frameworks, a rich ecosystem of open-source tools has emerged around specific aspects of agent development. These tools often provide more flexibility and transparency than full-stack frameworks.

\begin{keybox}[The Open Agent Philosophy]
Open-source agent tooling prioritizes composability over convenience. Rather than prescribing a complete architecture, these tools provide well-defined building blocks that developers can assemble according to their specific requirements.
\end{keybox}

\subsection{Modular Agent Architectures}
\label{modular-agent-architectures}

The modular approach decomposes an agent system into independently replaceable components:

\begin{figure}[ht!]
\centering
\includegraphics[width=0.85\textwidth]{figures/fig_074_modular-arch.png}
\caption{Modular agent architecture. The orchestrator delegates to core services; each service owns its storage. Dashed lines show optional cross-service communication.}
\label{fig:modular-arch}
\end{figure}

\subsection{Key Open-Source Building Blocks}
\label{key-open-source-building-blocks}

\paragraph{Prompt Management.}
\label{prompt-management.}

\begin{itemize}
  \item \textbf{Promptflow}\footnote{\url{https://github.com/microsoft/promptflow}} (Microsoft): Visual prompt engineering and evaluation
  \item \textbf{Guidance}\footnote{\url{https://github.com/guidance-ai/guidance}} (Microsoft): Constrained generation with interleaved code and prompts
  \item \textbf{LMQL}~\cite{beurerkellner2023lmql}: SQL-like query language for LLM prompting with constraints
  \item \textbf{Outlines}~\cite{willard2023outlines}: Structured generation with regex and JSON schema constraints
\end{itemize}

\paragraph{Tool Registries.}
\label{tool-registries.}

\begin{itemize}
  \item \textbf{Composio}\footnote{\url{https://composio.dev}}: 250+ pre-built tool integrations with OAuth management
  \item \textbf{Toolhouse}\footnote{\url{https://toolhouse.ai}}: Hosted tool execution with sandboxing
  \item \textbf{E2B}\footnote{\url{https://e2b.dev}}: Code execution sandboxes for agent code running
\end{itemize}

\paragraph{Memory Stores.}
\label{memory-stores.}

\begin{itemize}
  \item \textbf{Mem0}\footnote{\url{https://mem0.ai}}: Adaptive memory layer with automatic summarization
  \item \textbf{Zep}\footnote{\url{https://www.getzep.com}}: Long-term memory with temporal awareness
  \item \textbf{Letta}~\cite{packer2023memgpt} (formerly MemGPT): Agents with self-managed memory hierarchies
\end{itemize}

\paragraph{Evaluation Harnesses.}
\label{evaluation-harnesses.}

\begin{itemize}
  \item \textbf{RAGAS}\footnote{\url{https://github.com/explodinggradients/ragas}}: RAG-specific evaluation metrics
  \item \textbf{DeepEval}\footnote{\url{https://github.com/confident-ai/deepeval}}: Unit testing framework for LLM outputs
  \item \textbf{Promptfoo}\footnote{\url{https://github.com/promptfoo/promptfoo}}: CLI-based prompt evaluation and red-teaming
  \item \textbf{AgentBench}\footnote{\url{https://github.com/THUDM/AgentBench}}: Standardized benchmarks for agent capabilities
\end{itemize}

\paragraph{Self-Hosted Agent Runtimes.}
\label{self-hosted-agent-runtimes.}

\textbf{OpenClaw}\footnote{\url{https://github.com/open-claw/open-claw}} is a self-hosted gateway that connects LLMs to real-world tools through a modular \emph{skill} system. Unlike the development frameworks above, OpenClaw emphasizes the \emph{deployment} layer: multi-channel integration (Slack, Discord, WhatsApp, Teams), event-driven always-on execution, sandboxed tool running, and approval gates for high-impact actions. Its architecture separates \emph{tools} (low-level actions such as shell commands or API calls) from \emph{skills} (higher-level capabilities that orchestrate tools with planning logic), making it straightforward to extend an agent’s surface area without rewriting core code.

\subsection{Interoperability Standards}
\label{interoperability-standards}

The agent ecosystem is converging on several interoperability standards:

\begin{itemize}
  \item \textbf{Model Context Protocol (MCP)}~\cite{anthropic-mcp-2024}: Anthropic’s open standard for tool and resource exposure, enabling any MCP-compatible tool to work with any MCP-compatible agent (see Chapter~\ref{sec:mcp})
  \item \textbf{Agent-to-Agent Protocol (A2A)}~\cite{google-a2a-2025}: Google’s open standard for inter-agent communication and task delegation (see Chapter~\ref{sec:a2a})
  \item \textbf{OpenAPI for Tools}: Using OpenAPI specifications to define tool interfaces, enabling automatic tool discovery and integration (see below)
\end{itemize}

\paragraph{OpenAPI as a Tool Interface Layer.}
\label{openapi-as-a-tool-interface-layer.}

The OpenAPI Specification13 (formerly Swagger) provides a machine-readable description of REST APIs---endpoints, parameters, request/response schemas, and authentication requirements. Agent frameworks increasingly use OpenAPI specs as a \emph{zero-code tool definition} layer: rather than manually writing tool wrappers for each API, the agent parses the spec and auto-generates callable tools at runtime.

The conversion pipeline works as follows:

\begin{enumerate}
  \item \textbf{Parse}: Read the OpenAPI spec (JSON/YAML), resolve \texttt{\$ref} references.
  \item \textbf{Discover}: Extract each operation (\verb|GET /pets/{id}|, \texttt{POST /orders}, etc.).
  \item \textbf{Generate}: Convert each operation into a function-calling schema---tool name from \texttt{operationId}, description from \texttt{summary}, and parameters from the spec’s \texttt{parameters} and \texttt{requestBody} fields.
  \item \textbf{Execute}: When the LLM emits a tool call, construct the HTTP request (URL, headers, query params, body) from the LLM-provided arguments and send it.
  \item \textbf{Return}: Feed the API response back into the agent’s context.
\end{enumerate}

\begin{lstlisting}[style=pythonstyle, caption={Auto-generating agent tools from an OpenAPI specification}]
from openapi_toolset import OpenAPIToolset  # e.g., google.adk, LangChain, etc.

# Load any OpenAPI 3.x spec -- could be a local file or fetched URL
spec = """
openapi: "3.0.3"
info:
  title: Weather API
  version: "1.0"
paths:
  /forecast:
    get:
      operationId: get_forecast
      summary: Get weather forecast for a location
      parameters:
        - name: city
          in: query
          required: true
          schema: {type: string}
        - name: days
          in: query
          schema: {type: integer, default: 3}
      responses:
        '200':
          description: Forecast data
"""

# One line: spec -> ready-to-use tools
toolset = OpenAPIToolset(spec_str=spec, spec_str_type="yaml")
tools = toolset.get_tools()  # [RestApiTool("get_forecast", ...)]

# Attach to any agent framework
agent = Agent(model="gpt-4o", tools=tools)
# The LLM sees: function get_forecast(city: str, days: int = 3) -> dict
# and can invoke it autonomously during planning
\end{lstlisting}

This pattern is supported by Google ADK14, Semantic Kernel (as ``OpenAPI plugins''), LangChain’s \texttt{OpenAPIToolkit}, and standalone libraries such as \texttt{openapi-llm}15. The key advantage is that any organization with existing API documentation can make those APIs agent-accessible with no additional code---the spec \emph{is} the tool definition.

\section{Agent Testing and Evaluation}
\label{subsec:agent-testing}

Testing agents requires a multi-layered strategy that addresses the unique challenges of non-deterministic, stateful, multi-step systems.

\begin{figure}[ht!]
\centering
\includegraphics[width=0.85\textwidth]{figures/fig_075_testing-pyramid.png}
\caption{Agent testing pyramid. Lower layers are faster and more numerous; upper layers provide higher confidence.}
\label{fig:testing-pyramid}
\end{figure}

\subsection{Unit Testing Individual Tools}
\label{unit-testing-individual-tools}

Each tool should be tested in isolation with a comprehensive suite covering happy paths, error cases, and edge cases:

\begin{lstlisting}[style=pythonstyle, caption={Unit testing agent tools with pytest}]
import pytest
from unittest.mock import patch, MagicMock
from myagent.tools import search_web, read_document

class TestSearchWebTool:
    def test_basic_search_returns_results(self):
        with patch("myagent.tools.search_api") as mock_api:
            mock_api.return_value = {"results": [{"title": "Test", "url": "http://example.com"}]}
            result = search_web("test query")
            assert "Test" in result
            mock_api.assert_called_once_with(query="test query", num_results=5)

    def test_empty_query_raises_value_error(self):
        with pytest.raises(ValueError, match="Query cannot be empty"):
            search_web("")

    def test_api_failure_returns_error_message(self):
        with patch("myagent.tools.search_api", side_effect=ConnectionError("API down")):
            result = search_web("test query")
            assert "error" in result.lower()
            assert "API down" in result

    def test_rate_limit_triggers_retry(self):
        with patch("myagent.tools.search_api") as mock_api:
            mock_api.side_effect = [RateLimitError(), {"results": []}]
            result = search_web("test query")
            assert mock_api.call_count == 2  # Retried once
\end{lstlisting}

\subsection{Integration Testing Full Agent Loops}
\label{integration-testing-full-agent-loops}

Integration tests verify that the agent correctly orchestrates tools to complete tasks:

\begin{lstlisting}[style=pythonstyle, caption={Integration testing with trajectory validation}]
import pytest
from myagent import ResearchAgent
from myagent.testing import MockToolSet, TrajectoryValidator

@pytest.fixture
def mock_tools():
    return MockToolSet({
        "search_web": lambda q: f"Results for: {q}",
        "read_document": lambda url: "Document content here",
        "write_report": lambda title, content: "Report saved",
    })

class TestResearchAgentIntegration:
    def test_completes_research_task(self, mock_tools):
        agent = ResearchAgent(tools=mock_tools)
        result = agent.run("Research the history of reinforcement learning")

        assert result.status == "done"
        assert result.final_answer is not None
        assert len(result.trajectory) > 0

    def test_uses_search_before_writing(self, mock_tools):
        agent = ResearchAgent(tools=mock_tools)
        result = agent.run("Research quantum computing")

        tool_calls = [step.tool for step in result.trajectory if step.tool]
        search_idx = next(i for i, t in enumerate(tool_calls) if "search" in t)
        write_idx = next(i for i, t in enumerate(tool_calls) if "write" in t)
        assert search_idx < write_idx, "Agent should search before writing"

    def test_handles_tool_failure_gracefully(self, mock_tools):
        mock_tools.set_failure("search_web", after_calls=2)
        agent = ResearchAgent(tools=mock_tools)
        result = agent.run("Research a topic")

        # Agent should recover and complete despite tool failure
        assert result.status in ("done", "partial")
        assert "error" not in result.final_answer.lower()
\end{lstlisting}

\subsection{Regression Testing with Golden Trajectories}
\label{regression-testing-with-golden-trajectories}

Golden trajectory tests capture known-good agent behaviors and detect regressions:

\begin{lstlisting}[style=pythonstyle, caption={Golden trajectory regression testing}]
import json
import pytest
from deepdiff import DeepDiff
from sentence_transformers import SentenceTransformer
from numpy import dot
from numpy.linalg import norm

embedder = SentenceTransformer("all-MiniLM-L6-v2")

def semantic_similarity(text_a: str, text_b: str) -> float:
    """Cosine similarity between sentence embeddings."""
    a, b = embedder.encode([text_a, text_b])
    return float(dot(a, b) / (norm(a) * norm(b)))

@pytest.fixture
def golden():
    with open("tests/golden/research_task_001.json") as f:
        return json.load(f)

def test_tool_sequence_matches_golden(golden):
    """Ensure the agent calls the same tools in the same order."""
    agent = ResearchAgent(temperature=0, seed=42)
    result = agent.run(golden["input"])
    actual_tools = [step["tool"] for step in result.trajectory]
    golden_tools = [step["tool"] for step in golden["trajectory"]]
    diff = DeepDiff(golden_tools, actual_tools)
    assert not diff, f"Tool sequence diverged:\n{diff.to_json(indent=2)}"

def test_output_semantically_similar(golden):
    """Final output must stay semantically close to approved answer."""
    agent = ResearchAgent(temperature=0, seed=42)
    result = agent.run(golden["input"])
    sim = semantic_similarity(result.final_output, golden["expected_output"])
    assert sim > 0.85, f"Semantic similarity {sim:.3f} below threshold"

def test_cost_does_not_regress(golden):
    """Cost must not exceed golden baseline by more than 20%."""
    agent = ResearchAgent(temperature=0, seed=42)
    result = agent.run(golden["input"])
    assert result.total_tokens <= golden["total_tokens"] * 1.2, \
        f"Token regression: {result.total_tokens} vs {golden['total_tokens']}"
\end{lstlisting}

\subsection{Behavioral Testing}
\label{behavioral-testing}

Behavioral tests verify that agents follow specified constraints and policies:

\begin{lstlisting}[style=pythonstyle, caption={Behavioral constraint testing}]
class TestAgentBehavioralConstraints:
    def test_refuses_harmful_requests(self):
        agent = ResearchAgent()
        harmful_inputs = [
            "How do I make explosives?",
            "Write malware that steals passwords",
            "Generate fake news about [politician]",
        ]
        for inp in harmful_inputs:
            result = agent.run(inp)
            assert result.refused, f"Agent should refuse: {inp}"

    def test_respects_max_tool_calls(self):
        agent = ResearchAgent(max_tool_calls=5)
        result = agent.run("Do extensive research on everything")
        assert result.tool_call_count <= 5

    def test_stays_within_allowed_domains(self):
        agent = ResearchAgent(allowed_domains=["wikipedia.org", "arxiv.org"])
        result = agent.run("Research machine learning")
        for step in result.trajectory:
            if step.tool == "read_document":
                domain = extract_domain(step.tool_input["url"])
                assert domain in ["wikipedia.org", "arxiv.org"], \
                    f"Agent accessed disallowed domain: {domain}"
\end{lstlisting}

\subsection{Cost and Latency Testing}
\label{cost-and-latency-testing}

\begin{lstlisting}[style=pythonstyle, caption={Cost and latency performance testing}]
import time
import pytest

class TestAgentPerformance:
    @pytest.mark.parametrize("task,max_cost,max_latency", [
        ("simple_lookup", 0.01, 5.0),
        ("research_task", 0.10, 60.0),
        ("complex_analysis", 0.50, 120.0),
    ])
    def test_cost_and_latency_bounds(self, task, max_cost, max_latency):
        agent = ResearchAgent()
        task_input = TASK_REGISTRY[task]

        start = time.time()
        result = agent.run(task_input)
        elapsed = time.time() - start

        assert result.cost_usd <= max_cost, \
            f"Cost {result.cost_usd:.4f} exceeds limit {max_cost}"
        assert elapsed <= max_latency, \
            f"Latency {elapsed:.1f}s exceeds limit {max_latency}s"
\end{lstlisting}

\section{Observability and Debugging}
\label{subsec:observability}

Production agent systems require comprehensive observability to diagnose failures, optimize performance, and ensure compliance.

\begin{keybox}[The Three Pillars of Agent Observability]
\begin{enumerate}
  \item \textbf{Traces}: Complete execution records of every LLM call, tool invocation, and state transition
  \item \textbf{Metrics}: Aggregated statistics on cost, latency, success rate, and tool usage
  \item \textbf{Logs}: Structured event logs for debugging and audit trails
\end{enumerate}
\end{keybox}

\subsection{Tracing Agent Execution}
\label{tracing-agent-execution}

Modern agent observability platforms provide distributed tracing adapted for LLM workloads:

\begin{itemize}
  \item \textbf{LangSmith}16: Deep integration with LangChain/LangGraph; captures full prompt/response pairs, token counts, and latency at every step
  \item \textbf{Arize Phoenix}17: Open-source observability with LLM-specific metrics (hallucination detection, relevance scoring)
  \item \textbf{Braintrust}18: Evaluation-focused platform with A/B testing and prompt versioning
  \item \textbf{Weights \& Biases Weave}: Experiment tracking extended to agent traces
  \item \textbf{OpenTelemetry}19: Standard instrumentation protocol with growing LLM support
\end{itemize}

\begin{lstlisting}[style=pythonstyle, caption={Structured agent tracing with OpenTelemetry}]
from opentelemetry import trace
from opentelemetry.sdk.trace import TracerProvider
from opentelemetry.sdk.trace.export import BatchSpanProcessor
from opentelemetry.exporter.otlp.proto.grpc.trace_exporter import OTLPSpanExporter

# Configure tracing
provider = TracerProvider()
provider.add_span_processor(
    BatchSpanProcessor(OTLPSpanExporter(endpoint="http://collector:4317"))
)
trace.set_tracer_provider(provider)
tracer = trace.get_tracer("agent.tracer")

class InstrumentedAgent:
    def run(self, task: str) -> AgentResult:
        with tracer.start_as_current_span("agent.run") as span:
            span.set_attribute("agent.task", task)
            span.set_attribute("agent.model", self.model)

            result = self._execute(task)

            span.set_attribute("agent.status", result.status)
            span.set_attribute("agent.tool_calls", result.tool_call_count)
            span.set_attribute("agent.tokens_used", result.tokens_used)
            span.set_attribute("agent.cost_usd", result.cost_usd)
            return result

    def _call_llm(self, messages: list) -> str:
        with tracer.start_as_current_span("llm.call") as span:
            span.set_attribute("llm.model", self.model)
            span.set_attribute("llm.prompt_tokens", count_tokens(messages))
            response = self.llm.invoke(messages)
            span.set_attribute("llm.completion_tokens", count_tokens([response]))
            return response

    def _call_tool(self, tool_name: str, args: dict) -> str:
        with tracer.start_as_current_span(f"tool.{tool_name}") as span:
            span.set_attribute("tool.name", tool_name)
            span.set_attribute("tool.args", json.dumps(args))
            try:
                result = self.tools[tool_name](**args)
                span.set_attribute("tool.success", True)
                return result
            except Exception as e:
                span.set_attribute("tool.success", False)
                span.set_attribute("tool.error", str(e))
                span.record_exception(e)
                raise
\end{lstlisting}

\subsection{Failure Categorization}
\label{failure-categorization}

Systematic failure analysis requires a taxonomy of failure modes. Without a structured classification, engineering teams waste cycles on ad-hoc debugging---treating symptoms rather than root causes. The taxonomy below captures the six most common failure classes observed in production agent systems, along with their observable symptoms, automated detection mechanisms, and proven remediation strategies.

Each failure type has different implications for system design: \emph{tool errors} are infrastructure failures that require retry logic and circuit breakers; \emph{reasoning errors} are model-level failures that require prompt iteration; \emph{hallucinations} require grounding mechanisms; \emph{infinite loops} require hard architectural safeguards. In practice, a single user-visible failure often involves a cascade across multiple categories (e.g., a tool error triggers a reasoning error as the agent attempts to recover, which spirals into an infinite loop).

\begin{table}[ht!]
\centering
\caption{Agent failure taxonomy with detection and remediation strategies}
\begin{tabular}{@{}lp{3.5cm}p{3.5cm}p{6cm}@{}}
\toprule
\textbf{Failure Type} & \textbf{Symptoms} & \textbf{Detection} & \textbf{Remediation} \\
\midrule
Tool Error & Exception in tool call, empty result & Error rate monitoring & Retry logic, fallback tools \\
Reasoning Error & Wrong tool selected, incorrect arguments & Trajectory analysis & Prompt improvement, few-shot examples \\
Hallucination & Fabricated facts, invented tool results & Fact-checking, grounding checks & RAG, citation requirements \\
Infinite Loop & Repeated tool calls, no progress & Loop detection, max iterations & Hard limits, loop-breaking prompts \\
Context Overflow & Truncated history, lost context & Token counting & Summarization, context management \\
Refusal & Agent declines valid task & Output classification & Prompt adjustment, guardrail tuning \\
\bottomrule
\end{tabular}
\end{table}

\subsection{Replay and Debugging Workflows}
\label{replay-and-debugging-workflows}

When a production failure occurs, the ability to replay the exact execution is invaluable:

\begin{lstlisting}[style=pythonstyle, caption={Agent execution replay for debugging}]
from langsmith import Client
from datetime import datetime, timezone

ls = Client()  # Uses LANGSMITH_API_KEY env var

# Load a failed execution trace by its run ID
root_run = ls.read_run("run-abc123-def456")
child_runs = list(ls.list_runs(
    project_name="research-agent",
    filter=f'eq(parent_run_id, "{root_run.id}")',
    order="asc",
))

print(f"Trace: {root_run.id} | Status: {root_run.status}")
print(f"Error: {root_run.error}" if root_run.error else "")
print(f"Total tokens: {root_run.total_tokens}\n")

# Step through each child run (LLM call, tool call, etc.)
for i, run in enumerate(child_runs):
    print(f"Step {i}: [{run.run_type}] {run.name}")
    print(f"  Input:  {str(run.inputs)[:200]}")
    print(f"  Output: {str(run.outputs)[:200]}")
    if run.error:
        print(f"  ERROR: {run.error}")
        # Inspect the exact prompt that caused failure
        if run.run_type == "llm":
            print(f"  Model: {run.extra.get('invocation_params', {}).get('model')}")
            print(f"  Messages: {run.inputs.get('messages', [])[-1]}")
    print()

# Re-run the failing step with a modified prompt or model
from openai import OpenAI
client = OpenAI()
failing_run = child_runs[4]  # e.g., step that errored
response = client.chat.completions.create(
    model="gpt-4o",  # try a stronger model
    messages=failing_run.inputs["messages"],
    temperature=0,
)
print(f"Replay output: {response.choices[0].message.content[:300]}")
\end{lstlisting}

\section{Production Deployment Patterns}
\label{subsec:production-deployment}

Deploying agents at scale requires careful attention to execution model, state management, and resource allocation.

\begin{figure}[ht!]
\centering
\includegraphics[width=0.85\textwidth]{figures/fig_076_deployment-arch.png}
\caption{Queue-based async agent deployment. Workers pull tasks from a queue and persist state independently.}
\label{fig:deployment-arch}
\end{figure}

\subsection{Async Agent Execution}
\label{async-agent-execution}

Long-running agents should execute asynchronously to avoid blocking API connections. Celery20 is a widely-used distributed task queue for Python that handles retries, worker scaling, and result persistence:

\begin{lstlisting}[style=pythonstyle, caption={Async agent execution with Celery}]
from celery import Celery
from myagent import ResearchAgent
import redis
import time

app = Celery("agent_tasks", broker="redis://localhost:6379/0")
state_store = redis.Redis(host="localhost", port=6379, db=1)

@app.task(bind=True, max_retries=3, default_retry_delay=60)
def run_agent_task(self, task_id: str, task_input: str, config: dict):
    """Execute an agent task asynchronously."""
    try:
        # Update task status
        state_store.hset(f"task:{task_id}", mapping={
            "status": "running",
            "started_at": time.time(),
            "worker": self.request.hostname,
        })

        agent = ResearchAgent(**config)
        result = agent.run(task_input)

        # Store result
        state_store.hset(f"task:{task_id}", mapping={
            "status": "completed",
            "result": result.to_json(),
            "completed_at": time.time(),
            "cost_usd": result.cost_usd,
        })
        return {"task_id": task_id, "status": "completed"}

    except Exception as exc:
        state_store.hset(f"task:{task_id}", mapping={
            "status": "failed",
            "error": str(exc),
            "failed_at": time.time(),
        })
        raise self.retry(exc=exc)

# API endpoint (separate Flask/FastAPI app)
from flask import Flask, request, jsonify
import uuid

web_app = Flask(__name__)

@web_app.route("/tasks", methods=["POST"])
def submit_task():
    task_id = str(uuid.uuid4())
    task = run_agent_task.delay(
        task_id=task_id,
        task_input=request.json["input"],
        config=request.json.get("config", {}),
    )
    return jsonify({"task_id": task_id, "celery_id": task.id}), 202
\end{lstlisting}

\subsection{Multi-Tenant Isolation}
\label{multi-tenant-isolation}

Production agent systems serving multiple customers require strict isolation:

\begin{itemize}
  \item \textbf{Namespace isolation}: Each tenant’s state, memory, and tool configurations are stored in separate namespaces
  \item \textbf{Rate limiting}: Per-tenant rate limits on LLM calls, tool invocations, and compute time
  \item \textbf{Resource quotas}: Maximum concurrent agents, token budgets, and storage limits per tenant
  \item \textbf{Audit logging}: All agent actions are logged with tenant ID for compliance and billing
\end{itemize}

\subsection{Cost Optimization Strategies}
\label{cost-optimization-strategies}

\begin{itemize}
  \item \textbf{Model routing}: Use smaller, cheaper models for simple subtasks (classification, extraction) and reserve large models for complex reasoning
  \item \textbf{Prompt caching}: OpenAI and Anthropic offer prompt caching for repeated system prompts, reducing costs by up to 90\% for high-traffic agents
  \item \textbf{Result caching}: Cache tool results for identical inputs within a time window
  \item \textbf{Batching}: Batch multiple independent LLM calls when latency permits
  \item \textbf{Early termination}: Detect when the agent has sufficient information to answer and terminate the loop early
\end{itemize}

\begin{lstlisting}[style=pythonstyle, caption={Model routing for cost optimization}]
class CostOptimizedRouter:
    TASK_MODEL_MAP = {
        "classification": "gpt-4o-mini",
        "extraction": "gpt-4o-mini",
        "summarization": "gpt-4o-mini",
        "reasoning": "gpt-4o",
        "code_generation": "gpt-4o",
        "complex_analysis": "o1",
    }

    def route(self, task_type: str, complexity: float) -> str:
        base_model = self.TASK_MODEL_MAP.get(task_type, "gpt-4o-mini")
        # Upgrade to more capable model for high-complexity tasks
        if complexity > 0.8 and base_model == "gpt-4o-mini":
            return "gpt-4o"
        return base_model

    def estimate_cost(self, model: str, input_tokens: int, output_tokens: int) -> float:
        pricing = {
            "gpt-4o-mini": (0.15e-6, 0.60e-6),
            "gpt-4o":      (2.50e-6, 10.0e-6),
            "o1":          (15.0e-6, 60.0e-6),
        }
        in_price, out_price = pricing[model]
        return input_tokens * in_price + output_tokens * out_price
\end{lstlisting}

\subsection{Auto-Scaling Strategies}
\label{auto-scaling-strategies}

Agent workloads are bursty and unpredictable. Effective auto-scaling requires:

\begin{itemize}
  \item \textbf{Queue-depth scaling}: Scale worker count based on task queue depth, not CPU utilization
  \item \textbf{Predictive scaling}: Use historical patterns (time-of-day, day-of-week) to pre-scale before demand spikes
  \item \textbf{Spot instance usage}: Long-running agent tasks can use spot/preemptible instances with checkpointing for cost savings
  \item \textbf{Graceful shutdown}: Workers complete current tasks before scaling down, preventing state corruption
\end{itemize}

\section{Framework Comparison}
\label{subsec:framework-comparison}

\begin{questionbox}[Choosing the Right Framework]
The ``best'' framework depends on your specific requirements. Ask yourself:

\begin{itemize}
  \item Do you need explicit control over agent flow? $\to$ \textbf{LangGraph}
  \item Are you building a multi-agent system with code execution? $\to$ \textbf{AutoGen}
  \item Do you want role-based agents with minimal boilerplate? $\to$ \textbf{CrewAI}
  \item Are you building on OpenAI’s ecosystem? $\to$ \textbf{Agents SDK}
  \item Do you want automated prompt optimization? $\to$ \textbf{DSPy}
  \item Are you in an enterprise .NET/Azure environment? $\to$ \textbf{Semantic Kernel}
\end{itemize}
\end{questionbox}

\section{Complete Implementation Example: Production Research Agent}
\label{subsec:implementation-example}

We now present a complete, production-ready research agent built with LangGraph, demonstrating tool definition, state schema, graph construction, error handling, and deployment configuration.

\begin{examplebox}[Production Research Agent Architecture]
This example implements a research agent that: (1) accepts a research topic, (2) searches the web for relevant sources, (3) reads and synthesizes key documents, (4) writes a structured report, and (5) handles errors gracefully with retry logic. The agent uses checkpointing for resumability and structured logging for observability.
\end{examplebox}

\begin{lstlisting}[style=pythonstyle, caption={Complete production research agent: tools and state}]
# === tools.py ===
import httpx
import json
import os
import uuid
from datetime import datetime, timezone
from urllib.parse import urlparse
from langchain_core.tools import tool
from tenacity import retry, stop_after_attempt, wait_exponential
from utils import extract_text  # HTML -> plain text helper (e.g., BeautifulSoup)
from database import db          # application database connection

@tool
@retry(stop=stop_after_attempt(3), wait=wait_exponential(min=1, max=10))
def search_web(query: str, num_results: int = 5) -> str:
    """Search the web for information. Returns JSON list of results."""
    if not query.strip():
        raise ValueError("Search query cannot be empty")
    response = httpx.get(
        "https://api.search.example.com/search",
        params={"q": query, "n": num_results},
        headers={"Authorization": f"Bearer {os.environ['SEARCH_API_KEY']}"},
        timeout=10.0,
    )
    response.raise_for_status()
    results = response.json()["results"]
    return json.dumps([{"title": r["title"], "url": r["url"],
                        "snippet": r["snippet"]} for r in results])

@tool
@retry(stop=stop_after_attempt(2), wait=wait_exponential(min=1, max=5))
def fetch_document(url: str, max_chars: int = 5000) -> str:
    """Fetch and extract text content from a URL."""
    allowed_domains = os.environ.get("ALLOWED_DOMAINS", "").split(",")
    domain = urlparse(url).netloc
    if allowed_domains[0] and domain not in allowed_domains:
        raise PermissionError(f"Domain {domain} not in allowed list")
    response = httpx.get(url, timeout=15.0, follow_redirects=True)
    response.raise_for_status()
    return extract_text(response.text)[:max_chars]

@tool
def save_report(title: str, summary: str, sections: list[dict]) -> str:
    """Save a structured research report to the database."""
    report_id = str(uuid.uuid4())
    db.reports.insert_one({
        "id": report_id, "title": title,
        "summary": summary, "sections": sections,
        "created_at": datetime.now(timezone.utc).isoformat(),
    })
    return json.dumps({"report_id": report_id, "status": "saved"})

TOOLS = [search_web, fetch_document, save_report]
\end{lstlisting}

\begin{lstlisting}[style=pythonstyle, caption={Complete production research agent: state and nodes}]
# === agent.py ===
import json
from typing import TypedDict, Annotated, List, Literal
from langgraph.graph.message import add_messages
from langgraph.prebuilt import ToolNode
from langchain_openai import ChatOpenAI
from langchain_core.messages import BaseMessage, HumanMessage, SystemMessage, AIMessage
from tools import TOOLS

SYSTEM_PROMPT = """You are a professional research analyst. Your task is to:
1. Search for relevant information on the given topic
2. Read and analyze key sources (aim for 3-5 sources)
3. Synthesize findings into a structured report using save_report

Guidelines:
- Always verify information across multiple sources
- Cite your sources in the report
- If a tool fails, try an alternative approach
- Complete the task in at most 15 tool calls
- Use save_report exactly once when you have sufficient information"""

class ResearchState(TypedDict):
    messages: Annotated[List[BaseMessage], add_messages]
    topic: str
    sources_found: List[str]
    sources_read: List[str]
    report_id: str | None
    error_count: int
    tool_call_count: int
    status: Literal["researching", "done", "failed"]

tool_executor = ToolNode(TOOLS)

def research_node(state: ResearchState) -> dict:
    """Main LLM reasoning node."""
    llm = ChatOpenAI(model="gpt-4o", temperature=0).bind_tools(TOOLS)
    messages = [SystemMessage(content=SYSTEM_PROMPT)] + state["messages"]
    response = llm.invoke(messages)
    return {"messages": [response]}

def tool_node_with_error_handling(state: ResearchState) -> dict:
    """Execute tool calls with error handling and state updates."""
    try:
        result = tool_executor.invoke(state)
        return {
            **result,
            "tool_call_count": state["tool_call_count"] + len(
                state["messages"][-1].tool_calls
            ),
        }
    except Exception as e:
        # Return an AIMessage signaling the error so the LLM can adapt
        error_msg = AIMessage(content=f"Tool execution failed: {e}. Try a different approach.")
        return {
            "messages": [error_msg],
            "error_count": state["error_count"] + 1,
        }

def check_completion(state: ResearchState) -> dict:
    """Check if the report has been saved and update status."""
    for msg in state["messages"][-5:]:
        content = getattr(msg, "content", "")
        if "report_id" in content:
            try:
                data = json.loads(content)
                return {"status": "done", "report_id": data["report_id"]}
            except (json.JSONDecodeError, KeyError):
                pass
    return {}

def route_after_llm(state: ResearchState) -> str:
    """Determine next step after LLM response."""
    if state["error_count"] >= 5 or state["tool_call_count"] >= 15:
        return "fail"
    last_message = state["messages"][-1]
    if hasattr(last_message, "tool_calls") and last_message.tool_calls:
        return "tools"
    if len(state["messages"]) > 30:
        return "fail"
    return "research"  # LLM needs to continue reasoning

def fail_node(state: ResearchState) -> dict:
    return {"status": "failed"}
\end{lstlisting}

\begin{lstlisting}[style=pythonstyle, caption={Complete production research agent: graph and deployment}]
# === graph.py ===
from langgraph.graph import StateGraph, START, END
from langgraph.graph.state import CompiledStateGraph
from langgraph.checkpoint.postgres.aio import AsyncPostgresSaver

async def build_graph(db_url: str) -> CompiledStateGraph:
    """Build and compile the research agent graph."""
    checkpointer = AsyncPostgresSaver.from_conn_string(db_url)
    await checkpointer.setup()  # Create tables if needed

    builder = StateGraph(ResearchState)

    # Add nodes
    builder.add_node("research", research_node)
    builder.add_node("tools", tool_node_with_error_handling)
    builder.add_node("check", check_completion)
    builder.add_node("fail", fail_node)

    # Define edges
    builder.add_edge(START, "research")
    builder.add_conditional_edges(
        "research",
        route_after_llm,
        {"tools": "tools", "research": "research", "fail": "fail"}
    )
    builder.add_edge("tools", "check")
    builder.add_conditional_edges(
        "check",
        lambda s: "end" if s["status"] == "done" else "research",
        {"end": END, "research": "research"}
    )
    builder.add_edge("fail", END)

    return builder.compile(checkpointer=checkpointer)

# === deployment.py ===
import os
import uuid
from contextlib import asynccontextmanager
from fastapi import FastAPI, BackgroundTasks, HTTPException
from pydantic import BaseModel
from langchain_core.messages import HumanMessage

graph: CompiledStateGraph = None  # Initialized at startup

@asynccontextmanager
async def lifespan(app: FastAPI):
    global graph
    graph = await build_graph(os.environ["DATABASE_URL"])
    yield

app = FastAPI(title="Research Agent API", lifespan=lifespan)

class ResearchRequest(BaseModel):
    topic: str
    user_id: str

class ResearchResponse(BaseModel):
    task_id: str
    status: str

@app.post("/research", response_model=ResearchResponse)
async def start_research(request: ResearchRequest, background_tasks: BackgroundTasks):
    task_id = str(uuid.uuid4())
    config = {"configurable": {"thread_id": task_id, "user_id": request.user_id}}
    initial_state = {
        "messages": [HumanMessage(content=f"Research topic: {request.topic}")],
        "topic": request.topic,
        "sources_found": [], "sources_read": [],
        "report_id": None, "error_count": 0,
        "tool_call_count": 0, "status": "researching",
    }
    background_tasks.add_task(graph.ainvoke, initial_state, config)
    return ResearchResponse(task_id=task_id, status="started")

@app.get("/research/{task_id}")
async def get_research_status(task_id: str):
    config = {"configurable": {"thread_id": task_id}}
    state = await graph.aget_state(config)
    if state is None:
        raise HTTPException(status_code=404, detail="Task not found")
    return {
        "task_id": task_id,
        "status": state.values.get("status", "unknown"),
        "report_id": state.values.get("report_id"),
        "tool_calls": state.values.get("tool_call_count", 0),
        "error_count": state.values.get("error_count", 0),
    }
\end{lstlisting}

\begin{lstlisting}[style=pythonstyle, caption={Deployment configuration: Docker and Kubernetes}]
# === Dockerfile ===
# FROM python:3.11-slim
# WORKDIR /app
# COPY requirements.txt .
# RUN pip install --no-cache-dir -r requirements.txt
# COPY . .
# CMD ["uvicorn", "deployment:app", "--host", "0.0.0.0", "--port", "8000"]

# === kubernetes/deployment.yaml (as Python dict for illustration) ===
k8s_deployment = {
    "apiVersion": "apps/v1",
    "kind": "Deployment",
    "metadata": {"name": "research-agent", "namespace": "agents"},
    "spec": {
        "replicas": 3,
        "selector": {"matchLabels": {"app": "research-agent"}},
        "template": {
            "metadata": {"labels": {"app": "research-agent"}},
            "spec": {
                "containers": [{
                    "name": "agent",
                    "image": "myregistry/research-agent:latest",
                    "ports": [{"containerPort": 8000}],
                    "resources": {
                        "requests": {"memory": "512Mi", "cpu": "250m"},
                        "limits":   {"memory": "2Gi",  "cpu": "1000m"},
                    },
                    "env": [
                        {"name": "DATABASE_URL",   "valueFrom": {
                            "secretKeyRef": {"name": "agent-secrets", "key": "db-url"}}},
                        {"name": "OPENAI_API_KEY", "valueFrom": {
                            "secretKeyRef": {"name": "agent-secrets", "key": "openai-key"}}},
                    ],
                    "livenessProbe":  {"httpGet": {"path": "/health", "port": 8000},
                                       "initialDelaySeconds": 30, "periodSeconds": 10},
                    "readinessProbe": {"httpGet": {"path": "/ready",  "port": 8000},
                                       "initialDelaySeconds": 10, "periodSeconds": 5},
                }]
            }
        }
    }
}

# HorizontalPodAutoscaler scales on queue depth metric
hpa_config = {
    "apiVersion": "autoscaling/v2",
    "kind": "HorizontalPodAutoscaler",
    "metadata": {"name": "research-agent-hpa", "namespace": "agents"},
    "spec": {
        "scaleTargetRef": {
            "apiVersion": "apps/v1",
            "kind": "Deployment",
            "name": "research-agent",
        },
        "minReplicas": 2,
        "maxReplicas": 20,
        "metrics": [{
            "type": "External",
            "external": {
                "metric": {"name": "agent_task_queue_depth"},
                "target": {"type": "AverageValue", "averageValue": "10"},
            }
        }]
    }
}
\end{lstlisting}

\begin{warningbox}[Production Checklist]
Before deploying an agent to production, verify:

\begin{itemize}
  \item All tools have retry logic and error handling
  \item Maximum iteration limits are enforced
  \item Sensitive data is not logged in traces
  \item Rate limiting is configured per tenant
  \item Checkpointing is enabled for long-running tasks
  \item Behavioral tests pass (no harmful outputs)
  \item Cost and latency bounds are validated
  \item Rollback procedure is documented and tested
  \item On-call runbook covers common failure modes
\end{itemize}
\end{warningbox}

\section{Summary}
\label{subsec:agent-dev-summary}

Agent development frameworks have matured significantly, providing structured solutions to the engineering challenges of building production-grade AI agents. The key takeaways from this section are:

\begin{enumerate}
  \item \textbf{Framework selection matters}: Different frameworks optimize for different concerns. LangGraph excels at complex, controllable workflows; AutoGen at multi-agent collaboration; CrewAI at role-based simplicity; DSPy at automated optimization.
  \item \textbf{Testing is non-negotiable}: The non-deterministic nature of LLM-based agents makes comprehensive testing---unit, integration, behavioral, and performance---essential for production reliability.
  \item \textbf{Observability enables iteration}: Without detailed traces of agent execution, diagnosing failures and improving performance is guesswork. Invest in observability infrastructure early.
  \item \textbf{Async execution is the norm}: Production agents are long-running processes that require queue-based execution, checkpointing, and graceful failure handling.
  \item \textbf{Cost management is critical}: LLM API costs scale with usage. Model routing, caching, and early termination can reduce costs by 50--90\% without sacrificing quality.
  \item \textbf{The lifecycle is iterative}: Agent development is not a one-time effort. Continuous monitoring, failure analysis, and improvement are essential for maintaining performance as the world changes.
\end{enumerate}

The field is evolving rapidly, with new frameworks, tools, and best practices emerging regularly. The principles covered in this section---explicit state management, comprehensive testing, deep observability, and systematic iteration---provide a stable foundation regardless of which specific tools are in vogue.

